% Generated draft: requires mathematical and source review.
\documentclass[11pt]{article}
\usepackage[margin=25mm]{geometry}
\usepackage{fontspec}
\setmainfont{Latin Modern Roman}
\usepackage{amsmath,amssymb,mathtools,unicode-math}
\setmathfont{Latin Modern Math}
\usepackage{enumitem,needspace,xurl,hyperref}
\setlength{\parindent}{0pt}
\setlength{\parskip}{5pt}
\setlength{\emergencystretch}{3em}
\newcommand{\N}{\mathbb{N}}
\newcommand{\Z}{\mathbb{Z}}
\newcommand{\Q}{\mathbb{Q}}
\newcommand{\R}{\mathbb{R}}
\newcommand{\C}{\mathbb{C}}
\newcommand{\F}{\mathbb{F}}
\newcommand{\A}{\mathbb{A}}
\newcommand{\PP}{\mathbb{P}}
\newcommand{\E}{\mathbb{E}}
\newcommand{\eps}{\varepsilon}
\newcommand{\dd}{\,\mathrm{d}}
\begin{document}
\section*{UnsolvedMath ID: 10600058}
\textbf{Source problem code:} catalogue entry
\subsection{Definitions and mathematical statement}
\emph{Construction problem.} Construct a group for graph-links analogous to the group in the source's cited reference. The citation is absent from the supplied statement, so the comparison group, its presentation, and the required analogy are not defined.

A graph-link is understood here only as an equivalence class of the graph diagrams under the intended graph-link Reidemeister moves; the excerpt supplies neither its exact diagram class nor the moves. A proposed construction must specify whether it assigns an invariant group \(G(L)\) to each graph-link \(L\), or constructs a single group encoding graph-link diagrams, moves, or related data. In the invariant-group interpretation it must at least prove
\[
 L\cong L'\ \Longrightarrow\ G(L)\cong G(L'),
\]
with any required marked or peripheral structure specified. In a common-group interpretation the encoding map and its compatibility with moves would instead have to be stated. Neither interpretation is selected by the source excerpt. No fundamental-group analogue, complete invariant, or group operation on graph-link classes is silently imposed.
\subsection{Short English statement}
Construct a group-theoretic model for graph-links analogous to an omitted reference. Recovering that reference is essential: the excerpt does not say whether the group is an invariant of each graph-link or a shared algebraic model of diagrams and moves.
\subsection{Sources}
\begin{itemize}
\item Local UnsolvedMath registry, permanent ID 10600058, source code catalogue entry, virtual-knot problem 58. No individual proposer is supplied. The stored statement and title both have a missing citation after ``the group from''.
\item \emph{Unsolved Problems in Virtual Knot Theory and Combinatorial Knot Theory}, revised survey, problem-list item 58. Supplied URL: \url{https://arxiv.org/abs/1409.2823}; recorded access 2026-07-29; extraction source-tex. The dedicated section contains 60 top-level enumerated items with compound items retained intact.
\item The supplied research summary suggests the ``Reidemeister Moves and Groups'' sense and gives \texttt{arXiv:1412.xxxx} (2014). The identifier is an incomplete placeholder and does not identify a verified comparison construction.
\item Supplied status open and OPEN-TRIAGE. The summary says the problem appears open and no literature was found on this specific graph-link group. This is an editorial search assessment. Current primary-literature status and redistribution terms for the publicly accessible author/publisher copy remain NEEDS\_REVIEW.
\end{itemize}
\end{document}
